\chapter{Arithmetik mit Gedächtnis: der Halbring}
\label{ch:arithmetik}

% Register: D6, S5, S6, S10-S16, V9. Quellen: Spez. 8, Anhang A-C.

Bis hierher haben wir Hori-Zahlen nur betrachtet und verschoben.
Jetzt wollen wir mit ihnen rechnen. Das klingt harmlos,
ist aber die Stelle, an der die Horimetrik ihre erste echte
Entscheidung verlangt: Wenn zwei Hori-Zahlen addiert werden, welchen
Horizont bekommt das Ergebnis?

Im gewöhnlichen Zahlensystem stellt sich diese Frage nicht -- dort hat
ein Ergebnis keinen Horizont, nur einen Wert. Hier dagegen ist der
Horizont Teil der Zahl. Eine Rechenvorschrift muss deshalb aus zwei
Teilen bestehen: einer Regel für den Wert und einer Regel für den
Horizont. Die Horizontregel nennen wir einen \emph{Selektor}.

\section{Der erste Fehlschlag und was er lehrt}

Der naheliegendste Selektor ist der konservative: Das Ergebnis erbt
den größten beteiligten Horizont, notfalls vergrößert um das, was der
Ergebniswert mindestens braucht. Für die Addition funktioniert das
tadellos. Für die Multiplikation scheitert es -- und zwar an der
unschuldigsten aller Zahlen.

\begin{sackgasse}
Multipliziert man konservativ, dann gilt im Allgemeinen
$(a\htimes b)\htimes c\ne a\htimes(b\htimes c)$, sobald eine Null
beteiligt ist. Die Null löscht den \emph{Wert} jedes Produkts, aber
der konservative Selektor schleppt die \emph{Horizonte} der
Zwischenschritte weiter mit. Je nachdem, wann die Null zuschlägt,
erinnert sich das Ergebnis an verschiedene Rechenwege -- gleicher
Wert, verschiedene Horizonte, keine Klammerfreiheit. Lehre: Ein
Horizont trägt Rechengeschichte, und Rechengeschichte verträgt sich
nicht automatisch mit den Rechengesetzen.
\end{sackgasse}

\begin{beispiel}
Konkret, mit der konservativen Multiplikation (Horizont = Maximum
der Beteiligten, notfalls vergrößert): Es sei $a=b=(2,5)$ -- die
Fünf mit Horizont 2 -- und $c=(1,0)$ die Null mit Horizont 1. Links
geklammert: $a\htimes b=(4,25)$, denn $25$ braucht Horizont 4;
danach $(4,25)\htimes c=(4,0)$. Rechts geklammert:
$b\htimes c=(2,0)$, danach $a\htimes(2,0)=(2,0)$. Beide Wege enden
beim Wert null -- aber einmal mit Horizont 4, einmal mit Horizont 2.
$(4,0)\ne(2,0)$: keine Assoziativität.
\end{beispiel}

\section{Der wertseitige Halbring}

Der Ausweg ist eine Arbeitsteilung: Die Addition darf sich erinnern,
die Multiplikation muss vergessen. (Eine Konvention vorweg, die für
das ganze Buch gilt: \emph{Gesetzestreu} nennen wir eine Arithmetik,
wenn beide Operationen assoziativ und kommutativ sind, die
Multiplikation über die Addition distribuiert und ein additiv
neutrales Element existiert. Eine multiplikative Eins und das
Schlucken durch die additive Null verlangen wir \emph{nicht}
grundsätzlich, sondern nur dort, wo sie ausdrücklich genannt
werden -- beides wird uns gleich als echte Unterscheidungsmerkmale
begegnen.)

\begin{satz}[wertseitiger Halbring]\label{satz:a3}
Auf dem Horiraum $\Hori=\{(n,x)\mid n\ge\muh(x)\}$ seien, in
Wertkoordinaten $(n,x)$,
\[
(n,x)\hplus(m,y)=\bigl(\max(n,m,\muh(x+y)),\,x+y\bigr),
\qquad
(n,x)\htimes(m,y)=\bigl(\muh(xy),\,xy\bigr).
\]
Dann sind $\hplus$ und $\htimes$ kommutativ und assoziativ,
$\htimes$ distribuiert über $\hplus$, $(0,0)$ ist neutral für
$\hplus$ und absorbierend für $\htimes$.
\end{satz}

Der Beweis ist erstaunlich kurz; sein einziger Kern ist die
Monotonie des Minimalhorizonts: mehr Wert braucht nie weniger Platz
(vollständig in Anhang~\ref{ch:anhang-beweise}).
Bemerkenswert ist, was diesem Halbring \emph{fehlt}: eine Eins. Es
gibt kein Element, das bei Multiplikation alles unverändert lässt.
Stattdessen gilt für die Hori-Zahl Eins -- also $(1,1)$, die Eins in
ihrem kleinsten Zuhause -- die Gleichung
\[
h\htimes(1,1)=\KV(h)
\qquad\text{mit}\quad
\KV(n,x)=(\muh(x),\,x).
\]
Multiplikation mit der Eins ist die
Wertkompaktifizierung.\footnote{Terminologische Warnung:
„Kompaktifizierung“ ist hier nicht topologisch gemeint. Formal
handelt es sich um idempotente Retraktionen auf den jeweiligen
Minimalhorizont -- wir behalten das anschauliche Wort bei und bitten
topologisch geschulte Leser um Nachsicht.} Die Eins
wirkt nicht als neutrales Element, sondern als \emph{Projektor}: Sie
drückt jede Hori-Zahl auf ihre kompakte Form. Die kompakten Elemente
bilden darunter einen eigenen kleinen Halbring, und dieser ist exakt
das gewöhnliche $\Nn$. Die vertrauten natürlichen Zahlen sitzen also
als „Eck“ im Horiraum, und der Weg dorthin ist keine Zusatzoperation,
sondern Multiplikation mit der Eins.

\begin{beispiel}
Rechnen wir einmal durch. $(3,7)\hplus(2,5)$: Wert $12$, Horizont
$\max(3,2,\muh(12))=\max(3,2,3)=3$, also $(3,12)$ -- der größte
beteiligte Horizont überlebt. $(3,7)\htimes(2,5)$: Wert $35$,
Horizont $\muh(35)=4$, also $(4,35)$ -- die Vergangenheit ist
vergessen. Und der Projektor: Die Sieben mit unnötig großem
Horizont, $(5,7)$, mal die Eins $(1,1)$ ergibt
$(\muh(7),7)=(3,7)$ -- die Multiplikation mit Eins hat aufgeräumt.
\end{beispiel}

\begin{oberflaeche}
Was ist hier passiert? Wir haben eine Arithmetik gefunden, in der
Zahlen sich beim Addieren an ihre Vergangenheit erinnern -- der
größte beteiligte Horizont überlebt -- und beim Multiplizieren neu
geboren werden. Und das gewöhnliche Rechnen, das wir seit der
Grundschule kennen, ist darin enthalten: als der Teil des Horiraums,
der keine Vergangenheit mit sich trägt.
\end{oberflaeche}

\section{Exzess: das Gedächtnis bekommt eine Koordinate}

Lange sah es so aus, als sei diese vergessende Multiplikation
alternativlos. Die Vermutung lag nahe, dass die Rechengesetze jede
andere Horizontregel zerstören. Sie ist falsch, und die Widerlegung
ist lehrreicher als ein Beweis es gewesen wäre.

Der Horizont jeder Hori-Zahl zerfällt kanonisch in zwei Anteile:
\[
n=\underbrace{\muh(x)}_{\text{Pflicht}}
 +\underbrace{\varepsilon}_{\text{Exzess}} .
\]
Der Pflichtanteil ist das Minimum, das der Wert braucht; der Exzess
ist freiwillig mitgeführtes Gedächtnis. Diese Zerlegung ist keine
Konvention, sondern durch die Fakultätsschalen eindeutig bestimmt --
und sie entkoppelt die Frage nach dem Selektor vollständig: Man darf
auf dem Exzess \emph{eigenständig} rechnen.

\begin{satz}[Produktprinzip]\label{satz:produktprinzip}
Jede kommutative Halbringstruktur auf den Exzessen liefert eine
gesetzestreue Hori-Arithmetik: Werte rechnen gewöhnlich,
Exzesse rechnen in der gewählten Struktur.
\end{satz}

Das Produktprinzip liefert damit eine ganze Familie neuer
Arithmetiken. Zusammen mit dem schon konstruierten vergessenden
Halbring aus Satz~\ref{satz:a3} -- der, wohlgemerkt, \emph{keine}
Produktkonstruktion ist: Sein Additions-Selektor vergleicht die
Horizonte selbst, nicht nur die Exzesse -- ergeben sich drei
besonders natürliche Arithmetiken:

\begin{center}
\begin{tabular}{lccc}
 & Gedächtnis bei $+$ & Gedächtnis bei $\times$ & Eins \\
\hline
vergessend (Satz~\ref{satz:a3}; kein Produkt) & Maximum (der Horizonte) & stirbt & keine \\
bewahrend (tropisch) & Maximum & addiert sich & $(1,1)$ \\
verstärkend (arithmetisch) & addiert sich & multipliziert sich & $(2,1)$ \\
\end{tabular}
\end{center}

\begin{beispiel}
Es sei $a=(5,7)$ -- Pflicht $\muh(7)=3$, Exzess $2$ -- und
$b=(3,5)$ -- Pflicht $2$, Exzess $1$. Das Produkt hat den Wert $35$
mit Pflicht $\muh(35)=4$. Die drei Arithmetiken unterscheiden sich
nur im Exzess des Ergebnisses: vergessend $(4,35)$ (Exzess $0$),
bewahrend $(4{+}2{+}1,35)=(7,35)$ (Exzesse addieren sich),
verstärkend $(4{+}2,35)=(6,35)$ (Exzesse multiplizieren sich:
$2\cdot1=2$).
\end{beispiel}

Die bewahrende Arithmetik hat eine Kuriosität: Ihre Null schluckt
nicht. Multipliziert man eine Hori-Zahl mit null, wird der Wert
ausgelöscht, aber das Gedächtnis überlebt. Die verstärkende dagegen
ist ein vollwertiger kommutativer Halbring mit Eins und schluckender
Null -- und ihre Multiplikation löscht das Gedächtnis nicht
pauschal, sondern verrechnet es arithmetisch. Ganz gedächtnistreu
ist freilich auch sie nicht: Ein Operand mit Exzess null absorbiert
das Gedächtnis des anderen, denn null mal irgendetwas bleibt null --
genau daher rührt ihre schluckende Null.

\begin{oberflaeche}
Die Frage „welchen Horizont bekommt das Ergebnis?“ hat sich in
etwas viel Schöneres verwandelt: „Was soll mit der Vergangenheit
geschehen, wenn gerechnet wird?“ Und die Antwort der Horimetrik ist:
Das darfst du wählen -- vergessen, bewahren oder verstärken -- und
jede Wahl ergibt eine in sich vollkommen widerspruchsfreie Welt.
\end{oberflaeche}

\section{Die Dualität: dieselbe Eins, zwei Projektoren}

Bis hierher rechnete alles auf der Wertseite. Aber Hori-Zahlen
tragen neben dem Wert auch ihre Gestalt (und, seit
Kapitel~\ref{ch:erweiterungen}, den Bruch -- von dem in
Abschnitt~\ref{sec:seitentreue} noch die Rede sein wird) -- und die Gestaltseite hat
ihre eigene Arithmetik. Man
kann statt der Werte die \emph{Strukturpolynome} addieren und
multiplizieren; das ist das ziffernweise Rechnen ohne Überträge, bei
dem etwa aus den Strukturen von $4$ und $5$ direkt die Struktur von
$20$ entsteht.

\begin{satz}[Dualitätssatz]\label{satz:dualitaet}
Der Horiraum trägt einen zweiten, strukturseitigen Halbring: Addition
konservativ, Multiplikation kompakt -- aber gemessen am
Strukturhorizont $\sigmah$ statt am Werthorizont $\muh$. Gerechnet
wird in Strukturkoordinaten $(n,P)$; über die Ziffern-Bijektion aus
Kapitel~\ref{ch:horizonte} bezeichnen $(n,x)$ und das Paar aus $n$
und der Gestalt von $x$ dasselbe Element. In ihm gilt
\[
h\boxtimes(1,1)=\KS(h)
\qquad\text{mit}\quad
\KS(n,P)=(\sigmah(P),\,P),
\]
und $\KS$ ist mit beiden Operationen verträglich. Das Eck dieses
Halbrings ist der Strukturpolynomraum.
\end{satz}

Die Pointe verdient einen eigenen Absatz. Es ist \emph{dieselbe}
Hori-Zahl $(1,1)$ -- die Eins --, die auf der Wertseite den Projektor
$\KV$ und auf der Strukturseite den Projektor $\KS$ trägt. Zwei
Multiplikationen, ein Element. Und das älteste Ergebnis dieses
Projekts -- dass die beiden Aufräum-Operationen nicht vertauschen;
der nächste Abschnitt macht das formal -- entpuppt sich als das
Nichtkommutieren dieser zwei Multiplikationen mit ein und derselben
Eins.

\section{Aufräumen: der Kommutator und das Tiefengesetz}
\label{sec:kommutator}

Weil auf ihnen ein ganzer Themenstrang dieses Buchs ruht, stellen
wir die beiden Projektoren jetzt nebeneinander (Register D6):
\[
\KV(n,x)=(\muh(x),\,x),
\qquad
\KS(n,P)=(\sigmah(P),\,P).
\]
$\KV$ räumt \emph{wertkompakt} auf: gleicher Wert, kleinster
Horizont -- die Gestalt wird dabei neu codiert. $\KS$ räumt
\emph{strukturkompakt} auf: gleiche Gestalt, kleinster Horizont --
der Wert wandert dabei mit dem Auswertungspunkt nach unten. Beide
sind idempotent, aber sie kommutieren nicht, und die Abweichung ist
messbar. Wir nennen
\[
\delta(h)=\hor\bigl(\KS\KV(h)\bigr)-\hor\bigl(\KV\KS(h)\bigr)
\]
den \emph{Kompaktifizierungs-Kommutator}; seine negativen Werte
heißen \emph{Tiefe}.

\begin{beispiel}
Der Hausgeist dieses Buchs: $h=(6,6)$, die Sechs in
Terminaldarstellung. $\KV$ zuerst: $(6,6)\mapsto(3,6)$; dort hat
die Sechs die Gestalt $X{+}2$ mit $\sigmah=2$, also
$\KS\KV(h)=(2,5)$ -- Horizont~$2$. $\KS$ zuerst: Die konstante
Gestalt $6$ ist schon strukturkompakt, $\KS(h)=h$, danach
$\KV(h)=(3,6)$ -- Horizont~$3$. Also $\delta(h)=2-3=-1$: Die
Reihenfolge des Aufräumens entscheidet, wo man landet -- und als
was.
\end{beispiel}

Nach oben wächst $\delta$ vergleichsweise schnell; der kleinste
Zeuge für $+2$ ist die $720$ als Gestalt $\ff X3$ im Horizont $9$
(die „einzelne Eins an der vierten Stelle“ aus dem letzten
Interludium). Nach unten dagegen ist $\delta$ fakultätisch träge,
und diese Trägheit gehorcht einem exakten Gesetz. Nenne
\[
g(m)=m-\min\{s\mid M_s(m{+}1)\ge m!\}
\]
die \emph{Schalenlücke} der Schale $S_m$: wie weit unter $m$ die
Kapazität einer Gestalt liegen darf, deren Wert trotzdem noch $S_m$
erreicht ($M_s$ ist die maximal gefüllte Gestalt aus
Kapitel~\ref{ch:strukturpolynome}).

\begin{satz}[Tiefengesetz, untere Seite; Register S33]
\label{satz:tiefengesetz}
Für jede Schale $m\ge3$ gilt
\[
\min\{\delta(h)\mid \muh(\val(h))=m\}=-g(m),
\]
und ein stets tauglicher Zeuge -- im Allgemeinen nicht der
einzige -- ist die Terminaldarstellung des Maximalpolynom-Wertes
$M_{m-g(m)}(m{+}1)$. Die Sprungstellen $m_g=\min\{m\mid g(m)\ge g\}$
bilden eine Zählfolge mit bewiesener Asymptotik
$m_g/(g{+}2)!\to1$ (Beweise in Anhang~\ref{ch:anhang-beweise}).
\end{satz}

Für Schale $3$ ist der Zeuge genau der Hausgeist von oben, und
$m_1=3$: Ab der dritten Schale ist Tiefe $-1$ möglich, ab der
fünfzehnten $-2$, ab der vierundachtzigsten $-3$. Die
\emph{beidseitigen Fixpunkte} schließlich sind die Elemente, denen
das Aufräumen nichts mehr anhaben kann -- fixiert von $\KV$
\emph{und} $\KS$ zugleich, also $n=\muh(x)=\sigmah(P)$. Beide
Begriffe, Sprungstellen wie Fixpunkte, werden in
Kapitel~\ref{ch:zaehlungen} gezählt.

Erwähnt sei noch das Feinkorn dieser Nicht-Kommutativität: Das von
$\KV$ und $\KS$ erzeugte Transformationsmonoid -- das
\emph{Kompaktifizierungsmonoid} -- ist endlich: genau sechs
Elemente, von denen alle außer $\KV\KS$ idempotent sind
(Register S21; Beweis in Anhang~\ref{ch:anhang-beweise}). Der
Zeuge mit dem kleinsten \emph{Wert} ist wieder der Hausgeist
$(6,6)$; in der kleinsten \emph{Welt} scheitert die Idempotenz von
$\KV\KS$ schon an $(5,25)$.

\section{Seitentreue}
\label{sec:seitentreue}

Kann man die beiden Seiten verweben -- etwa wertseitig addieren und
strukturseitig multiplizieren? Gesetzestreu heißt dabei, wie
eingangs vereinbart: beide Operationen kommutativ und assoziativ,
die Multiplikation distribuiert, eine additive Null existiert. Der
systematische Test aller sechzehn Kombinationen -- Wert oder
Gestalt als Payload, konservativ oder kompakt als Selektor, für
jede der beiden Operationen -- gibt dann eine klare Antwort: Genau
zwei \emph{dieser sechzehn} sind gesetzestreu, der wertseitige
Halbring und sein strukturseitiger Spiegel. (Die Produktfamilie aus
Satz~\ref{satz:produktprinzip} steht außerhalb dieses Rasters --
ihre Selektoren rechnen auf Exzessen -- und ist durchweg
gesetzestreu, bleibt aber ganz auf der Wertseite.) Alle acht gemischten
Kandidaten sterben an der Distributivität; die übrigen sechs
seitenreinen scheitern an der Neutralität (kompakte Addition
vergisst den Horizont des Operanden) oder an der Assoziativität
(konservative Multiplikation, das Nullproblem vom Kapitelanfang).

\begin{beispiel}
Der paradigmatische Zeuge, Ziffer für Ziffer. ($\oplus$ und
$\otimes$ bezeichnen in diesem Beispiel die gemischte
Kandidaten-Arithmetik, nicht die Operationen aus
Satz~\ref{satz:a3}; \emph{wertgetragen} heißt: die Operation
verrechnet die Werte, \emph{strukturgetragen}: die Gestalten.)
Es seien $a=b=(1,1)$ und
$c=(2,5)$, dessen Gestalt $X+2$ ist (Ziffern $12$). Addiert man
wertgetragen, ist $b\oplus c=(3,6)$; die Sechs hat in $H_3$
zufällig ebenfalls die Gestalt $X+2$ (Ziffern $012$). Multipliziert
man nun strukturgetragen mit $a$ (Gestalt $1$), bleibt die Gestalt
$X+2$ und wird kompakt bei Horizont 2 ausgewertet: Ergebnis
$(2,5)$ -- \emph{der Wert ist von 6 auf 5 gefallen}, weil die
Strukturmultiplikation nur die Gestalt kennt. Rechts dagegen:
$a\otimes b=(1,1)$, $a\otimes c=(2,5)$, wertgetragen addiert
$(3,6)$. Links $(2,5)$, rechts $(3,6)$: Die Distributivität ist
tot, und man sieht genau, woran -- die eine Seite rechnet mit dem
Wert ehrlich weiter, die andere hat ihn beim Gestaltwechsel
verloren.
\end{beispiel}

\begin{satz}[Seitentreue-Satz]\label{satz:seitentreue}
Jede gesetzestreue Arithmetik \emph{der hier betrachteten
payload-getra\-genen Form} -- die Operationen tragen Wert, Gestalt
oder Bruch, die Horizontregeln sind beliebig -- ist seitenrein:
Addition und Multiplikation operieren beide auf der Wertseite oder
beide auf der Strukturseite.
\end{satz}

Die Einschränkung auf payload-getragene Operationen ist wesentlich:
Auf einer abzählbaren Menge lässt sich per Transport jede beliebige
Halbringstruktur definieren -- der Satz spricht über die
Arithmetiken, die Wert, Gestalt oder Bruch \emph{respektieren},
und über diese vollständig.

Zur genauen Reichweite: Der Beweis braucht je nach Fall sogar
weniger als die volle Gesetzestreue -- ein Mischfall stirbt an der
Links-Distributivität allein, der andere an der beidseitigen
Distributivität (die zu jedem, auch nichtkommutativen, Halbring
gehört): Die Rechts-Distributivität spiegelt die Kettenkonstruktion
und ersetzt die Kommutativität. Auch nichtkommutative Mischformen
sind damit ausgeschlossen; denkbar bleibt allein ein nur
\emph{einseitig} distributives Mischsystem -- eine Kuriosität
außerhalb jeder Halbring-Axiomatik.

Der Beweis verdient es, in seiner Grundidee erzählt zu werden, denn
er hat eine Pointe (vollständig in Anhang~\ref{ch:anhang-beweise}). Zunächst sah es schlecht aus: Einzelne
Rechnungen, die eine gemischte Arithmetik erfüllen müsste, sind
tatsächlich lösbar -- teils über Pell-Gleichungen, also genau die
Sorte Zahlentheorie, die gerne Ausnahmen produziert. Ein direkter
Widerspruch schien unerreichbar. Die entscheidende Waffe ist
verblüffend schlicht: die \emph{Konstanten}. Multipliziert man eine
per Addition aufgebaute Kette von Zahlen mit immer größeren
konstanten Strukturen, dann treibt die Ziffernschranke die
Auswertungshorizonte ins Unendliche -- und ein nicht-konstantes
Polynom kann denselben Wert nicht an beliebig vielen Stellen
annehmen. Das zwingt die Kettenglieder, konstant zu sein. Konstanten
aber sterben am nächsten Multiplikator, weil auch dort die
Ziffernschranke den Horizont linear anwachsen lässt, während die
Distributivität lineares Wertwachstum verlangt. Beide Mischformen
zerbrechen an derselben Stelle.

Der Satz hat inzwischen sogar eine dritte Dimension: Auch die
Bruch-Identität aus Kapitel~\ref{ch:erweiterungen} lässt sich weder
mit der Wert- noch mit der Gestaltseite verweben, und eine eigene
Arithmetik trägt sie nicht (ihre Addition kann die Eins
überschreiten). Die Seitentreue ist eine Trilogie: drei Identitäten,
zwei Rechenwelten, keine Mischung -- Beweise in
Anhang~\ref{ch:anhang-beweise}.

Die Pointe: Die Quelle der Starrheit ist die Ziffernbedingung
$0\le a_i\le i$ -- die allererste, unschuldigste Regel dieses Buchs.
Die Regel, mit der alles begann, ist genau die Regel, die verhindert,
dass Wert und Gestalt jemals in einer Arithmetik verschmelzen. Starr
sind nicht die Horizontregeln; starr sind die Identitäten.

\begin{oberflaeche}
Eine Hori-Zahl hat drei Identitäten: was sie wert ist, wie sie
aussieht, und welcher Bruch sie ist. Dieses Kapitel hat gezeigt,
dass die ersten beiden vollwertige, spiegelbildliche Rechenwelten
tragen, dass ein und dieselbe Eins in beiden Welten als Türsteher
arbeitet -- und, als bewiesener Satz: dass niemand in zwei Welten
gleichzeitig rechnen kann, während die dritte gar keine eigene
trägt. Man muss sich entscheiden. Das ist der tiefste Satz, den die
Horimetrik zu bieten hat: Konsistenz hat einen Preis, und der Preis
ist eine Identität.
\end{oberflaeche}
